% 第三章：自由振动分析公式推导
% 包含：混合离散动力学方程、能量方程、弱形式

本章针对 Mindlin 中厚板的自由振动问题，基于第二章建立的运动学方程，推导混合离散格式下的动力学控制方程与广义特征值问题。首先推导板的动能变分与应变能变分，通过哈密顿原理建立自由振动的虚功方程；其次进行离散化，得到材料刚度矩阵与质量矩阵的各子块表达式；最后讨论混合离散格式的稳定性条件。

\section{混合离散动力学方程与特征值问题}

考虑 Mindlin 中厚板的自由振动问题，求解目标为结构的固有频率 $\omega$ 和固有振型。自由振动对应的广义特征值问题为：
\begin{equation}
(\mathbf{K} - \omega^2 \mathbf{M})\mathbf{D} = \mathbf{0}
\end{equation}
式中，$\mathbf{K}$ 为弹性刚度矩阵，$\mathbf{M}$ 为质量矩阵，$\mathbf{D}$ 为节点自由度向量。

\section{振动分析能量方程}

由第二章建立的 Mindlin 板位移场、应变-位移关系及材料本构可知，本章直接在此基础上进行动力分析。首先推导板的动能变分与应变能变分，再通过哈密顿原理建立自由振动的虚功方程。板的动能变分为：
\begin{equation}
\delta T = \int_V \rho \ddot{u}_i \delta u_i\, dV
\end{equation}
将位移场代入，利用中面位置 $x_3=0$ 的对称性以及密度均匀假设，定义惯性系数
\begin{equation}
I_0 = \int_{-h/2}^{h/2} \rho\, dx_3 = \rho h, \qquad I_2 = \int_{-h/2}^{h/2} \rho x_3^2\, dx_3 = \frac{\rho h^3}{12}
\end{equation}
由于中面对称性，$\int_{-h/2}^{h/2} \rho x_3 dx_3 = 0$，可得动能变分的中面积分形式：
\begin{equation}
\delta T = \int_\Omega \left(I_0 \ddot{w}\,\delta w + I_2 \ddot{\psi}_\alpha \delta\psi_\alpha\right) d\Omega
\end{equation}

板的应变能变分为：
\begin{equation}
\delta U = \int_V \sigma_{ij} \delta\varepsilon_{ij}\, dV
\end{equation}
内虚功密度 $\sigma_{ij}\delta\varepsilon_{ij}$ 包含面内分量与横向分量（$\sigma_{33}=0$ 不计）：
\begin{equation}
\sigma_{ij}\delta\varepsilon_{ij} = \sigma_{\alpha\beta}\delta\varepsilon_{\alpha\beta} + 2\sigma_{\alpha 3}\delta\varepsilon_{\alpha 3}
\end{equation}
利用 Mindlin 位移场对应的应变变分 $\delta\varepsilon_{\alpha\beta} = x_3\delta\kappa_{\alpha\beta}$，$\delta\varepsilon_{\alpha 3} = \frac{1}{2}\delta\gamma_\alpha$，代入并对厚度积分，定义弯矩合力与剪力合力
\begin{equation}
M_{\alpha\beta} \equiv \int_{-h/2}^{h/2} x_3 \sigma_{\alpha\beta}\, dx_3 = \frac{h^3}{12} C_{\alpha\beta\gamma\delta}\kappa_{\gamma\delta}, \qquad Q_\alpha \equiv \int_{-h/2}^{h/2} \sigma_{\alpha 3}\, dx_3 = \kappa G h \gamma_\alpha
\end{equation}
由此，内虚功可写为
\begin{equation}
\delta U = \int_\Omega \left(M_{\alpha\beta}\delta\kappa_{\alpha\beta} + Q_\alpha \delta\gamma_\alpha\right) d\Omega
\end{equation}

根据哈密顿原理，系统在任意时间区间 $[t_1, t_2]$ 内的动能变分与应变能变分满足：
\begin{equation}
\int_{t_1}^{t_2} (\delta T - \delta U)\, dt = 0,
\end{equation}
对任意虚位移 $\delta u_i$，可得弹性动力学的虚功方程：
\begin{equation}
\int_V (\sigma_{ij}\delta\varepsilon_{ij} + \rho \ddot{u}_i \delta u_i)\, dV = 0.
\end{equation}
板的体积积分可分解为中面积分与厚度积分：$\int_V dV = \int_\Omega \int_{-h/2}^{h/2} dx_3\, d\Omega$，其中 $\Omega$ 为板的中面。将内虚功与惯性项代入，得到 Mindlin 板自由振动的虚功方程：
\begin{equation}
\int_\Omega \left[M_{\alpha\beta}\delta\kappa_{\alpha\beta} + Q_\alpha \delta\gamma_\alpha - I_0 \ddot{w}\,\delta w - I_2 \ddot{\psi}_\alpha \delta \psi_\alpha\right] d\Omega = 0.
\end{equation}

\section{弱形式}

对上述虚功方程各变量取变分并令其为零，将离散场代入后得到广义特征值方程。位移场和转角场的形函数插值为：
\begin{equation}
w(\mathbf{x}) = N_I w_I, \qquad \varphi_\alpha(\mathbf{x}) = N_I \varphi_{\alpha I}
\end{equation}
对空间坐标求偏导，微分算子直接作用在形函数上：$w_{,\alpha} = N_{I,\alpha}w_I$，$\varphi_{\alpha,\beta} = N_{I,\beta}\varphi_\alpha$。虚位移做同阶离散近似。

材料刚度矩阵各子块为
\begin{equation}
K_{IJ}^{ww} = \int_\Omega \kappa G h\, N_{I,\alpha} N_{J,\alpha}\, d\Omega, \qquad K_{I,J\gamma}^{w\varphi} = -\int_\Omega \kappa G h\, N_{I,\gamma} N_J\, d\Omega
\end{equation}
\begin{equation}
K_{I\alpha,J\gamma}^{\varphi\varphi} = \int_\Omega \left(N_{I,\beta} D_{\alpha\beta\gamma\eta} N_{J,\eta} + \kappa G h\, N_I N_J \delta_{\alpha\gamma}\right) d\Omega
\end{equation}

质量矩阵展开后的对称矩阵形式为
\begin{equation}
\mathbf{M} = \begin{bmatrix}
M^{ww} & \mathbf{0} & \mathbf{0} \\
\mathbf{0} & M^{\varphi_1\varphi_1} & \mathbf{0} \\
\mathbf{0} & \mathbf{0} & M^{\varphi_2\varphi_2}
\end{bmatrix}
\end{equation}
其中横向挠度对应的质量子块为：
\begin{equation}
M_{IJ}^{ww} = \int_\Omega I_0 N_I N_J\, d\Omega = \int_\Omega \rho h\, N_I N_J\, d\Omega
\end{equation}
转角自由度对应的质量子块为：
\begin{equation}
M_{I\alpha,J\beta}^{\varphi\varphi} = \int_\Omega I_2 \delta_{\alpha\beta} N_I N_J\, d\Omega = \int_\Omega \frac{\rho h^3}{12} \delta_{\alpha\beta} N_I N_J\, d\Omega
\end{equation}

由于线性问题中面内与弯曲不耦合，全局特征值方程呈现块对角结构：
\begin{equation}
\begin{bmatrix}
\mathbf{K}_{ww} & \mathbf{K}_{w\varphi} \\
\mathbf{K}_{\varphi w} & \mathbf{K}_{\varphi\varphi}
\end{bmatrix}
\begin{bmatrix}
\mathbf{W} \\
\boldsymbol{\Psi}
\end{bmatrix}
= \omega^2
\begin{bmatrix}
\mathbf{M}_{ww} & \mathbf{0} \\
\mathbf{0} & \mathbf{M}_{\varphi\varphi}
\end{bmatrix}
\begin{bmatrix}
\mathbf{W} \\
\boldsymbol{\Psi}
\end{bmatrix}
\end{equation}
此即 Mindlin 板自由振动的全局离散特征值方程。求解广义特征值问题 $(\mathbf{K} - \omega^2\mathbf{M})\mathbf{D} = \mathbf{0}$，可得到自然频率 $\omega$ 和对应模态。

格式稳定性方面，振动问题中，质量矩阵 $\mathbf{M}$ 为对称正定。为保证自由振动特征值问题的适定性，混合离散格式需满足 LBB 稳定性条件。本文采用的剪切应力场插值阶数经约束比分析验证，可保证 inf-sup 常数有正下界，从而在薄板极限下避免剪切自锁并维持刚度矩阵的正定性。当板厚趋近于零时，最低阶频率收敛于 Kirchhoff 薄板理论的解析解。
